%% ma: L12-B11-0006
%% lop: 12
%% chuong: 4
%% bai: 11
%% dang: 12.11.3
%% loai: TN4
%% muc: TH
%% dap_an: "D"
%% nguon: "(NBV)-ĐỀ SỐ 8-MỖI NGÀY 1 ĐỀ THI 2025.pdf (D08, MỖI NGÀY 1 ĐỀ - Đề số 8), Phần I Câu 6"
%% kien_thuc: "Bài 11 (nguyên hàm của hàm lượng giác), Lớp 11 Bài 2 (công thức lượng giác: sin²+cos²=1, hiệu hai bình phương)"
%% trang_thai: chinh-thuc
%% ghi_chu: "Đáp án gốc D đúng. Cùng ý tưởng biến đổi lượng giác trước khi lấy nguyên hàm với L12-B11-0003 (dạng 12.11.3)."
\begin{ex}
Nguyên hàm của hàm số $f(x)=\cos^4\dfrac x2-\sin^4\dfrac x2$ là
\choice{$-\cos x+C$}{$\cos x+C$}{$-\sin x+C$}{\True $\sin x+C$}
\loigiai{$f(x)=\left(\cos^2\dfrac x2-\sin^2\dfrac x2\right)\left(\cos^2\dfrac x2+\sin^2\dfrac x2\right)=\cos x$.
Do đó
\[\int f(x)\,\mathrm dx=\int\cos x\,\mathrm dx=\sin x+C.\]}
\end{ex}
